%======================================================================
% 附录I：学习策略与刻意练习手册
%======================================================================
% 章标题 \appchapter{附录I：...} 由主文件 plasma_physics_textbook_v2.tex 发出
% （检查套件按主文件判定附录字母连续性），本文件只提供正文。

前面八个附录回答的是"等离子体物理是什么"，这一个附录回答的是"怎么把它学进脑子里"。

本书是自学教材：没有老师点名，没有考试倒逼，也没有人替你判断"到底懂没懂"。自学最大的风险通常不是内容太难，而是\textbf{用低效的方法把时间花光}——花了三个月，书翻旧了，合上书却推不出德拜长度。

本附录把学习科学中经过实验检验的做法，翻译成可以直接执行的动作，并逐条指出它在本书里的落点。

\begin{insight}{本附录与前言的分工}{}
前言里的"三类自学人群路线图"回答\textbf{先读什么、读多快、跳过什么}；本附录回答\textbf{读完之后怎么巩固、隔多久复习、怎么证明自己真的会了}。两者合起来才是一条完整的学习闭环：前言负责排课，本附录负责验收。
\end{insight}


\section{三个把"读懂了"误当"学会了"的模式}

\begin{warning}{流畅性错觉}{}
重新阅读、划线、抄公式，都会让材料\textbf{看起来}越来越熟悉。这种熟悉感来自"再认"（看到答案能认出来），而不是"回忆"（不给提示能写出来）。考试与科研需要的是后者，而前者会稳定地高估自己的掌握程度。
\end{warning}

\textbf{误区一：重读与划线。}
把一章读三遍，比读一遍再合书复述一遍，效果差得多：前者提高的是对排版的熟悉度，后者才是在建立提取路径。
诊断方法很简单：读完一节后立刻合上书，在白纸上写出这一节的\textbf{核心方程、适用条件、一个反例}。写不出来，说明刚才那一遍只是"看过了"。

\textbf{误区二：集中突击。}
连续 6 小时攻同一章，短期感觉很爽；两周后回看，遗忘得比"每天 1 小时、分 6 天"更彻底。
同样总时长下，把学习\textbf{分散}到多天，长期保持率显著更高——这条结论在跨度几十年的元分析中反复出现（见 I.7 文献表）。

\textbf{误区三：同类题目连做。}
把 10 道双流不稳定性题目连着做完，正确率会虚高，因为每道题都在重复同一个套路；一旦题目混在一起、需要自己判断"该用哪个方法"，正确率立刻掉下来。
更有效的做法是\textbf{交错练习}：把第 6 章的波、第 7 章的动理学、第 9 章的不稳定性题目打乱混做，逼迫自己在解题前先做一次"方法选择"。

本书中，第 9 章补充材料的不稳定性分类树、附录 F 的 23 条高频易错点，都是为这种"先判断再计算"的训练准备的。


\section{五个可执行动作}

\subsection{动作一：合书复述（检索练习）}

\textbf{怎么做：}读完一节，立刻合书，用 3 分钟在白纸上写出：核心方程 $\to$ 每个符号的物理含义 $\to$ 推导的关键一步 $\to$ 成立条件 $\to$ 一个会被它证伪的反例。写完再翻书核对，\textbf{用红笔标出漏掉的那一条}。

\textbf{在本书的落点：}每章的"本章速查卡"（补充材料末尾的 `formulabox`）就是现成的核对清单；每章末尾的"自学检查清单"（`checklistbox`）可直接当复述提纲。

\textbf{常见走样：}变成"抄笔记"。抄写时眼睛一直盯着书，提取路径没有被使用，效果接近于零。

\subsection{动作二：间隔重复（1 / 3 / 7 / 21 天）}

\textbf{怎么做：}同一节内容，在学完后的第 1、3、7、21 天各做一次 2--3 分钟的合书复述。四次都过关，可以降为每月一次。

\textbf{为什么是这四个点：}遗忘在最初几天最快。复习的价值不在于"再看一遍"，而在于\textbf{在快要忘掉的时候把它取回来}——取回越费力（在成功的前提下），记忆越牢。

\textbf{在本书的落点：}把每章速查卡抄成 12 张卡片（一章一张），在卡片背面只写标题与日期，正面留白。复习时只看日期与标题，凭记忆默写符号表。

\subsection{动作三：交错练习}

\textbf{怎么做：}每周留出一次"混合练习"：从本周学习内容与之前 2--3 章中各抽 3 道题，打乱顺序限时完成。重点是\textbf{在动笔前用一句话写出"我为什么选这个方法"}。

\textbf{在本书的落点：}附录 K 的八道跨教材例题覆盖德拜屏蔽、磁镜、漂移、Alfvén 波、Landau 阻尼、Spitzer 电阻率、鞘层与重联，天然适合交错使用；把它们与各章例题混排出题效果最好。

\subsection{动作四：自我解释}

\textbf{怎么做：}每推完一步问自己三句话——"这一步用了什么假设？""如果这个假设不成立会怎样？""能不能用一句不含公式的话说出来？"把答案写在一旁。

\textbf{在本书的落点：}本书大量使用\textbf{推导后紧跟假设讨论}的写法（例如第 4 章流体闭合、第 7 章 Landau 阻尼的围道选择）。读到这些段落时不要跳过，它们正是自我解释的范本。

\subsection{动作五：例题读完合书重做}

\textbf{怎么做：}例题不要"读懂"就算过。读完题目后合上书，从空白纸开始独立做一遍；做不出来再看提示，然后\textbf{隔一天再做一遍}。同一道题做三遍（读后、隔日、隔周）比做三道同类新题更有价值。

\textbf{在本书的落点：}第 1--12 章的补充材料各含 5 道经典例题，加上附录 K 的 8 道，共 86 道；建议至少把其中 30 道按"三遍法"完整做过。

\begin{miniquiz}
\textbf{自测：}你的复习计划是"每周末把本周内容重读一遍"。请指出这个计划的两个缺陷，并改写成符合动作一与动作二的版本。（参考：重读不产生提取；一周一次的点太少且过晚。）
\end{miniquiz}


\section{元认知与错题本}

\textbf{错题本不是抄题本。}每道错题只记三栏：

\begin{enumerate}
\item \textbf{错在哪一步}——精确到"把 $\grad p$ 当成了标量""忘了 $\ee$ 的符号""单位换成了 CGS 但公式是 SI"；
\item \textbf{根因归类}——概念混淆（回第几章哪一节）／数学工具（第 1 章）／量纲与量级／审题遗漏；
\item \textbf{下次触发条件}——写成一个可执行的检查动作，例如"凡是出现 $\grad$ 与 $\vect{E}$ 同时出现的式子，先确认量纲"。
\end{enumerate}

三栏里最值钱的是第三栏。它把一次具体错误升级为一条\textbf{可复用的检查动作}，这正是"刻意练习"里"即时反馈 + 针对性重复"的落点：练习必须有明确目标、即时反馈与超出舒适区的难度，单纯重复已经会做的事不叫刻意练习。

\begin{checklistbox}{每周复盘清单}{app-i-weekly}
\begin{itemize}
\item 本周学过的每一节，都做过至少一次\textbf{合书复述}吗？
\item 错题本这周新增了几条？每条都写了"下次触发条件"吗？
\item 有没有做过一次\textbf{混排限时}练习（动作三）？
\item 有没有哪道例题，我在\textbf{不看答案}的情况下完整推出过？
\item 下周要复习的旧章节，是哪几章（按 1/3/7/21 天节奏）？
\end{itemize}
\end{checklistbox}


\section{量级估算与量纲检验}

物理直觉不是玄学，它由一批\textbf{记得住的典型数值}和一套\textbf{三十秒能做完的检查}组成。附录 K 的例题大多可以在两三行内估出结果，请把它们当作"直觉素材"背下来。

\subsection{三个必须记住的量级（氘等离子体）}

\begin{itemize}
\item \textbf{德拜长度}：$T=10\ \mathrm{eV}$、$n=10^{18}\ \mathrm{m^{-3}}$ 时，$\lambda_D\approx 23.5\ \mathrm{\mu m}$，对应 $N_D\approx 5.4\times 10^{4}$ 个粒子——远大于 1，所以流体图像在这里是合法的。
\item \textbf{回旋半径}：$\mathrm{D^+}$ 在 $B=5\ \mathrm{T}$ 下，$T=10\ \mathrm{eV}$ 时 $\rho_L\approx 0.13\ \mathrm{mm}$；$T=1\ \mathrm{keV}$ 时 $\approx 1.3\ \mathrm{mm}$；$T=10\ \mathrm{keV}$ 时 $\approx 4.1\ \mathrm{mm}$。$\rho_L\propto\sqrt{T}$，温度涨 1000 倍只要半径涨约 32 倍。
\item \textbf{Alfvén 速度}：$\mathrm{D^+}$、$n=10^{20}\ \mathrm{m^{-3}}$、$B=5\ \mathrm{T}$ 时 $v_A\approx 7.7\times 10^{6}\ \mathrm{m/s}$，约合光速的 2.6\%；磁场扰动穿过 1 m 只需约 $1.3\times 10^{-7}\ \mathrm{s}$。
\end{itemize}

\subsection{量纲检验三条}

\begin{enumerate}
\item \textbf{逐项同量纲}：方程两边、以及每一项，量纲必须一致。把 $\omega$、$k$、$v_A$、$\lambda_D$ 的量纲列成一张小表，比每次现推更快。
\item \textbf{无量纲组合}：物理结论通常只依赖无量纲量（$k\lambda_D$、$n\lambda_D^3$、$\beta$、$\rho_L/L$、$\omega/\omega_{ci}$）。先问"这个问题的无量纲参数是什么"，往往比直接积分更快看清物理。
\item \textbf{极限行为}：令参数取 $0$ 或 $\infty$ 看结果是否退化成已知情形。Landau 阻尼率在 $k\lambda_D\to 0$ 时应趋于零（$\propto\exp[-1/(2k^2\lambda_D^2)]$ 指数小），若推导结果在这极限下发散，一定错了。
\end{enumerate}

\begin{miniquiz}
\textbf{估算练习：}不查书，估算 $B=3\ \mathrm{T}$、$T_i=5\ \mathrm{keV}$ 的氘离子回旋频率与回旋半径。提示：$\omega_{ci}=qB/m$，$\rho_L=\sqrt{2T/m}/\omega_{ci}$（量级上用 $\sqrt{T/m}$ 即可）。答案应为 $\omega_{ci}\approx 1.4\times 10^{8}\ \mathrm{s^{-1}}$、$\rho_L$ 量级几毫米——如果你算出厘米或微米，请回第 3 章重做量纲。
\end{miniquiz}


\section{十二周节奏模板}

下面的模板假定\textbf{每天 60 分钟、每周 6 天、共 12 周}完成第一遍通读。它与前言的"三类人群路线图"配合使用：前言决定你读哪些章节，本表决定每天这 60 分钟怎么切、周末交什么作业。

\small
\begin{longtable}{@{}p{0.9cm}p{3.1cm}p{6.8cm}p{3.8cm}@{}}
\caption{十二周自学节奏模板（每天 60 分钟 $\times$ 6 天）}\label{tab:app-i-12weeks}\\
\toprule
周次 & 主攻内容 & 每天 60 分钟怎么切 & 周末产出 \\
\midrule
\endfirsthead
\toprule
周次 & 主攻内容 & 每天 60 分钟怎么切 & 周末产出 \\
\midrule
\endhead
\midrule
\multicolumn{4}{r}{\small（续下页）}\\
\endfoot
\bottomrule
\endlastfoot
1 & 第 1 章 数学与统计力学预备知识 & 15 min 复习昨日复述 $\to$ 25 min 读新节 $\to$ 15 min 例题重做 $\to$ 5 min 写速查卡 & 第 1 章清单自评；量纲表一张 \\
2 & 第 2 章 什么是等离子体 & 同上，例题时间加到 20 min & 德拜屏蔽三问（判据/量级/反例） \\
3 & 第 3 章 单粒子运动 & 20 min 读 $\to$ 20 min 画轨道图 $\to$ 20 min 估算 $\rho_L$、$\omega_c$ & 三类漂移方向默写；磁镜例题重做 \\
4 & 第 4 章 流体描述基础 & 25 min 推导 $\to$ 20 min 自我解释假设 $\to$ 15 min 速查卡 & 双流体方程闭合条件复述 \\
5 & 第 5 章 磁流体力学 & 25 min 读 $\to$ 20 min 例题 $\to$ 15 min 交错练习（混入第 4 章） & MHD 波三支速度量级估算 \\
6 & 第 6 章 等离子体中的波 & 20 min 读 $\to$ 25 min 画色散曲线 $\to$ 15 min 极限检验 & 冷等离子体色散关系默写 \\
7 & 第 7 章 动理学理论 & 25 min 读 $\to$ 20 min 分布函数练习 $\to$ 15 min 复述 Landau 围道 & Landau 阻尼率量级估算 \\
8 & 第 8 章 碰撞与输运 & 20 min 读 $\to$ 25 min 输运系数估算 $\to$ 15 min 错题本 & Spitzer 电阻率与扩散时间估算 \\
9 & 第 9 章 等离子体不稳定性 & 20 min 读 $\to$ 25 min 判据练习 $\to$ 15 min 分类树默画 & 不稳定性分类树；三种不稳定性判据对比 \\
10 & 第 10 章 约束与受控核聚变 & 25 min 读 $\to$ 20 min 劳森判据估算 $\to$ 15 min 交错练习 & 劳森判据三参数扫描手算 \\
11 & 第 11 章 低温工业等离子体 & 20 min 读 $\to$ 25 min 鞘层估算 $\to$ 15 min 复述玻姆判据 & 鞘层厚度与离子通量估算 \\
12 & 第 12 章 诊断 + 全局复习 & 30 min 诊断原理 $\to$ 30 min 附录 F/G/H 串读 & 全书 23 条易错点自查；附录 H 符号表抽测 \\
\end{longtable}
\normalsize

\textbf{进度落后怎么办：}宁可推迟一周，也不要把"合书复述"和"例题重做"砍掉——那两项才是模板的全部价值所在。可以先缓读第 11、12 章等相对独立的章节，但第 1--9 章的概念依赖链不能跳。


\section{自测题四型}

自测题的质量决定复习的质量。下面四型覆盖了从记忆到迁移的完整梯度，每型给出本书中的一个示例与判分标准。

\begin{enumerate}
\item \textbf{判据型——"什么时候可以用这个方法？"}\\
例：什么时候可以把等离子体当流体处理，什么时候必须用动理学？\\
判分：必须同时说出\textbf{判据}（$N_D\gg 1$、$k\lambda_D\ll 1$、碰撞足够频繁使分布接近 Maxwell 分布）与\textbf{失效应}（无碰撞、波--粒子共振时流体图像失效）。

\item \textbf{量级型——"答案大概是多少？"}\\
例：$B=5\ \mathrm{T}$、$T=10\ \mathrm{eV}$ 的氘离子回旋半径是多少？\\
判分：答案应落在 $0.1$--$0.2\ \mathrm{mm}$；数量级对即算过，但必须能说出 $\rho_L\propto\sqrt{T}/B$ 的来源。

\item \textbf{反例型——"这句话反着说对不对？"}\\
例："温度升高，碰撞频率一定升高。"请给出反例。\\
判分：指出库仑碰撞频率 $\nu\propto nT^{-3/2}$ 随温度\textbf{下降}，而与中性粒子的碰撞才随温度升高；能区分两种碰撞机制即算过。

\item \textbf{迁移型——"把方法搬到新问题上。"}\\
例：磁镜损失锥的几何推理（$\sin^2\theta_c=1/R$）能否搬到环形约束装置的俘获粒子分析上？\\
判分：能指出哪一部分可以搬（绝热不变量与速度空间判据）、哪一部分不能（磁场几何与 $B_{\max}/B_{\min}$ 的构造方式不同）。
\end{enumerate}

建议每章学完后，自己按四型各出一道题并交换给同伴（或自我隔周作答）。\textbf{能出题，才说明你看清了这一章的骨架。}


\section{出处与延伸阅读}

本附录的主张来自学习科学文献与经典解题方法论，列在下面以便读者自行核对。物理教材的对照关系见附录 J，分级书单见附录 E。

\normalsize
{\small
\begin{longtable}{@{}p{4.4cm}p{5.2cm}p{5.2cm}@{}}
\caption{本附录引用的学习科学文献与出处}\label{tab:app-i-refs}\\
\toprule
文献 & 出处 & 本附录中的用途 \\
\midrule
\endfirsthead
\toprule
文献 & 出处 & 本附录中的用途 \\
\midrule
\endhead
\midrule
\multicolumn{3}{r}{\small（续下页）}\\
\endfoot
\bottomrule
\endlastfoot
Dunlosky, Rawson, Marsh, Nathan \& Willingham (2013) & \textit{Perspectives on Psychological Science in the Public Interest} 14(1), 4--58；DOI: 10.1177/1529100612453266 & 十种学习技术的效用评级：检索练习与分散练习为高效用 \\
Roediger \& Karpicke (2006) & \textit{Psychological Science} 17(3), 249--255；DOI: 10.1111/j.1467-9280.2006.01693.x & 检索练习优于重复阅读（动作一） \\
Karpicke \& Blunt (2011) & \textit{Science} 331(6018), 772--775；DOI: 10.1126/science.1199327 & 检索练习优于概念图；强调"提取"本身 \\
Cepeda, Pashler, Vul, Wixted \& Rohrer (2006) & \textit{Psychological Bulletin} 132(3), 354--380；DOI: 10.1037/0033-2909.132.3.354 & 分散练习的元分析：间隔优于集中（动作二） \\
Rohrer \& Taylor (2007) & \textit{Instructional Science} 35(6), 481--498；DOI: 10.1007/s11251-007-9015-8 & 交错练习优于分块练习（动作三） \\
Chi, Bassok, Lewis, Reimann \& Glaser (1989) & \textit{Cognitive Science} 13(2), 145--182；DOI: 10.1207/s15516709cog1302\_1 & 自我解释与例题学习中的"好学生"行为（动作四） \\
Sweller \& Cooper (1985) & \textit{Cognition and Instruction} 2(1), 59--89；DOI: 10.1207/s1532690xci0201\_3 & 例题效应：先学例题再独立解题（动作五） \\
Ericsson, Krampe \& Tesch-R\"omer (1993) & \textit{Psychological Review} 100(3), 363--406；DOI: 10.1037/0033-295x.100.3.363 & 刻意练习三要素：明确目标、即时反馈、超出舒适区 \\
Bjork (1994) & 载 \textit{Metacognition: Knowing about Knowing}, MIT Press, 185--205 & "合意困难"：让练习变难反而学得更牢 \\
Polya (1945) & \textit{How to Solve It}, Princeton University Press & 解题四步法与本附录"元认知三问" \\
\end{longtable}
}
\normalsize

\textbf{一句总结：}把"我读过了"改成"我\textbf{不看任何提示}能写出核心方程、说出适用条件、给出一个反例"，自学的效率会发生数量级的变化。
